\chapter{CONCLUSION AND FUTURE SCOPE}
The proposed AI-based multi-label food spoilage detection system demonstrates the application of Deep Learning, Computer Vision, and Explainable AI techniques for automated food quality monitoring. The system was developed to classify fruits and vegetables into four spoilage stages: Fresh, Early Spoilage, Moderate Spoilage, and Severe Spoilage using image-based analysis. Multiple deep learning architectures including CustomCNN, MobileNetV2, ResNet50, and EfficientNetB0 were trained and evaluated using the same dataset and pre-processing pipeline for comparative analysis.

The experimental results indicate that MobileNetV2 achieved the best overall performance among all evaluated models with improved accuracy, F1-score, and ROC-AUC values. The lightweight architecture and efficient feature extraction capability of MobileNetV2 enabled better classification performance and faster inference compared to deeper architectures such as ResNet50 and EfficientNetB0. The confusion matrix and ROC analysis demonstrated that the proposed system was capable of performing effective multi-label spoilage classification, although intermediate spoilage stages remained more challenging because of overlapping visual characteristics.

The integration of Explainable AI using Grad-CAM improved model transparency and interpretability by highlighting image regions influencing predictions. The generated heatmaps demonstrated that the models focused on meaningful spoilage indicators such as discoloration, mold formation, bruises, and texture degradation during classification. The Streamlit-based web application further demonstrated the practical applicability of the proposed system for real-time food quality monitoring and prediction.

Although the IoT integration component is currently under development, the proposed framework successfully establishes the foundation for combining image-based deep learning and environmental sensing for food spoilage monitoring. The project demonstrates the potential of AI-driven systems for reducing food waste, improving food safety, and enabling intelligent food quality assessment in practical applications such as smart kitchens, grocery stores, warehouses, and inventory management systems.

\section[Future Scope]{\textbf{Future Scope}}
The future scope of the project includes complete integration of IoT-based environmental sensing using MQ-135 gas sensors and DHT11/DHT22 temperature-humidity sensors for real-time spoilage monitoring. Sensor fusion techniques can further improve prediction reliability by combining environmental conditions with image-based analysis. Additional food categories such as dairy products, packaged foods, and meat products can also be included to expand the applicability of the system.

Future work may also involve the use of larger and more diverse datasets to improve model generalization capability and classification accuracy for intermediate spoilage stages. Advanced architectures such as Vision Transformers (ViTs), multimodal learning models, and ensemble methods can be explored for improved performance. Cloud deployment, mobile application integration, and edge-AI implementation using embedded devices can further enhance the real-world usability of the proposed system.

\section[Learning Outcomes of the Project]{\textbf{Learning Outcomes of the Project}}
The project provided practical understanding of Deep Learning, transfer learning, image pre-processing, data augmentation, Explainable AI, and web application deployment techniques. The implementation process improved knowledge related to CNN architectures, transfer learning models, Grad-CAM visualization, and performance evaluation methods such as confusion matrix and ROC analysis.

The project also provided exposure to AI-based food quality monitoring systems and the integration of computer vision techniques with IoT-based monitoring frameworks. Practical implementation using TensorFlow, Keras, OpenCV, Streamlit, and Google Colab enhanced technical understanding of real-world AI application development workflows.
